\documentclass[10pt,a4paper]{amsart}
\usepackage[margin=19mm]{geometry}
\usepackage{amsmath,amssymb,mathtools}
\usepackage[scr=rsfs,cal=euler]{mathalfa}
\usepackage{xfrac}
\usepackage[unicode,hidelinks,bookmarks=false]{hyperref}
\newcommand{\ie}{i.e.\ }
\newcommand{\cf}{cf.\ }
\newcommand{\resp}{respectively\ }
\newcommand{\Subsection}{\subsection}
\providecommand{\nobreakdash}{\nobreak\hbox{-}}
\setlength{\emergencystretch}{2em}
\multlinegap0pt
\usepackage{times}
\usepackage{url}
\usepackage{booktabs}
\usepackage{nicefrac}
\newtheorem{theorem}{Theorem}
\newtheorem{corollary}[theorem]{Corollary}
\newtheorem{lemma}[theorem]{Lemma}
\newtheorem{definition}[theorem]{Definition}
\begin{document}
\title[Structure of the Circular-Dyadic Convolution Error]{Structure of the Circular-Dyadic Convolution Error}
\author{Ben Fauber and Alireza Moradzadeh; NVIDIA; Ben Fauber; Alireza Moradzadeh}
\date{}
\maketitle
\noindent\emph{Source and licence.} Structure of the Circular-Dyadic Convolution Error. Version arXiv:2607.15293v1 (CC BY 4.0). This material is distributed under the Creative Commons Attribution 4.0 International licence, http://creativecommons.org/licenses/by/4.0/, as recorded in the arXiv OAI licence metadata for this exact version and shown on the abstract page https://arxiv.org/abs/2607.15293v1. Full rights notice: licenses/arxiv-2607.15293-CC-BY-4.0.txt.\par\smallskip
\noindent\textit{Editorial scope.} Counted unit: the original \texttt{theorem} environment \texttt{thm:null-space} at source lines 210--226 together with its complete proof at lines 419--419; the delivered document keeps source lines 77--449 so that every referenced label and every citation resolves inside it. Native editable source is the paper's own concatenated TeX; the AMS wrapper is editorial formatting only. No publisher class, upstream script or unrelated artwork is redistributed.
\par\smallskip\noindent\textit{Reference note.} This article ships no compiled bibliography (biblatex); the entries delivered here are composed from the author's own BibTeX (.bib) fields, with no field invented and no entry dropped.
\begin{equation}
   \label{eq:circ-conv}
   y_{\mathrm{C}}[n] = \sum_{k=0}^{N-1} u[k]\,v\bigl[(n-k)\bmod N\bigr]
\end{equation}

\subsection{Hadamard Transform}
\label{sec:hadamard-transform}

The Hadamard (Walsh--Hadamard) transform is a generalized Fourier transform~\cite{walsh1923closed, shanks1969ITCmp.100.457S, beauchamp1975walsh}. The $(k,n)$-th entry of its associated $N \times N$ operator $\mathcal{H}_m$ is:
 
 \begin{equation}                                      
 \label{eq:hadamard-entry}
   (\mathcal{H}_{m})_{k,n}={\frac {1}{2^{m/2}}}(-1)^{\sum _{r} k_{r} n_{r}}
\end{equation}

\noindent where the sum runs over bit positions $r = 0, 1, \ldots, m-1$ and $k_r, n_r \in \{0,1\}$ are the $r$-th bits of $k$ and $n$. Since $\mathcal{H}_m$ is involutory ($\mathcal{H}_{m}^2 = I_{2^m}$), $\mathcal{H}_m^{-1} = \mathcal{H}_m$. It can be computed in $O(N\log N)$ operations \cite{beauchamp1975walsh, fino1976unified}, matching FFT complexity, using only sign flips and no complex arithmetic.

\medskip
\noindent\textbf{Remark ($N = 2^m$ requirement).} The restriction to powers of two is intrinsic: the recursive butterfly structure of $\mathcal{H}_m$ requires halving at each stage, and the XOR-based index arithmetic of equation~\eqref{eq:hadamard-entry} is defined over $m$-bit integers. Extensions to non-power-of-two lengths are outside the scope of this work.

\subsection{Hadamard Transform Convolution}
\label{sec:hadamard-conv}

$\mathcal{H}_m$ supports its own convolution theorem \cite{doi:10.1049/el:19730172, beauchamp1975walsh, uvsakova2002walsh}, computing a \emph{dyadic convolution}\footnote{The term \emph{dyadic} follows the Walsh--Hadamard convolution literature~\cite{doi:10.1049/el:19730172, beauchamp1975walsh}, where it denotes XOR‑indexed convolution over $(\mathbb{Z}_2)^m$. This usage is distinct from the wavelet‑theoretic (Haar) meaning of dyadic, which refers to multi-resolution analysis over dyadic intervals~\cite{ahmed1975orthogonal}.} that is not interchangeable with the circular convolution of the DFT under any output reordering (Section~\ref{sec:alg-structures} and Theorem~\ref{thm:zero-error}). The dyadic convolution is defined as follows:

\begin{equation}
      \label{eq:dyadic-conv}
       y_{\mathrm{D}}[n] = \sum_{k=0}^{N-1} u[k]\,v[n \oplus k]
\end{equation}

\noindent where $\oplus$ denotes XOR, replacing the modular indexing $(n-k)\bmod N$ of equation~\eqref{eq:circ-conv}. The XOR indexing arises from the binary structure of equation~\eqref{eq:hadamard-entry}: the sign $(-1)^{\sum_r k_r n_r}$ decomposes independently at each bit position, and XOR is the index arithmetic that preserves this structure. The following convolution identity holds:\footnote{The $\sqrt{N}$ factor is a normalization artifact of the unitary convention $\mathcal{H}_m^2 = I$; the unnormalized matrix $H_m$ (entries $\pm1$, $H_m^{-1} = H_m/N$) gives the equivalent form $y_{\mathrm{D}} = \tfrac{1}{N}H_m\{H_m\{u\} \odot H_m\{v\}\}$, directly analogous to equation~\eqref{eq:circ-conv}.}

\begin{equation}
      \label{eq:hadamard-thm}
      y_{\mathrm{D}} = \sqrt{N}\;\mathcal{H}_m^{-1}\!\left\{\mathcal{H}_m\{u\} \odot \mathcal{H}_m\{v\}\right\}
\end{equation}

\subsection{Difference in Algebraic Structures}
\label{sec:alg-structures}

Although $\mathcal{H}_m$ is equivalent to a multidimensional DFT at the transform level~\cite{shanks1969ITCmp.100.457S, yarlagadda1997hadamard}, this equivalence does not extend to convolution over $\mathbb{Z}_N$. Both operations are group convolutions; the term \emph{circular} throughout this paper refers specifically to convolution over the cyclic group $\mathbb{Z}_N$ (modular indexing), not to group convolution in general. The DFT diagonalizes circular convolution, whereas $\mathcal{H}_m$ diagonalizes dyadic convolution over the binary group $(\mathbb{Z}_2)^m$ via XOR indexing \cite{Terras_1999}. These groups are not isomorphic for $m \geq 2$: $\mathbb{Z}_N$ contains an element of order $2^m$, while every element of $(\mathbb{Z}_2)^m$ has order at most $2$. Consequently, the two convolutions are algebraically distinct. In the framework of P\"uschel and Moura~\cite{PueschelMoura2008_ASP_1DSpace, PueschelMoura2008_ASP_Foundation1DTime}, this distinction reflects that the DFT and $\mathcal{H}_m$ are Fourier transforms for different, non‑isomorphic signal models; the structured error produced by substituting one for the other is the focus of Section~\ref{sec:results} onward.

\section{Results}
\label{sec:results}

Throughout Section~\ref{sec:results}, let $N = 2^m$ and $u, v \in \mathbb{R}^N$; we define the \emph{error signal}:

\begin{equation}
    e[n] \;\triangleq\; y_{\mathrm{C}}[n] \;-\; y_{\mathrm{D}}[n]
\end{equation}

\noindent where $y_{\mathrm{C}}$ and $y_{\mathrm{D}}$ are the circular and dyadic convolutions of equations~\eqref{eq:circ-conv} and~\eqref{eq:dyadic-conv}. Let $\mathbf{C}^u$ and $\mathbf{H}^u$ denote the $N\times N$ matrices with $(\mathbf{C}^u)_{n,k} = u[(n-k)\bmod N]$ and $(\mathbf{H}^u)_{n,k} = u[n\oplus k]$, so that $\mathbf{e} = (\mathbf{C}^u - \mathbf{H}^u)v$. The full error vector is written $\mathbf{e}$ (bold) to distinguish it from the scalar $e[n]$; $\mathbf{b}_j$ denotes the $j$-th standard basis vector. Signal vectors $u$, $v$, $y_{\mathrm{C}}$, $y_{\mathrm{D}}$ are written unbolded throughout.

\subsection{Foundational Zero-Error Conditions}
\label{sec:foundational}

The first result characterizes exact error cancellation: we identify the output and input positions where both convolutions agree for all $u, v \in \mathbb{R}^N$.

\begin{theorem}[Universal zero-error output positions]\label{thm:zero-error}
For all $N = 2^m$ with $m \geq 2$, there are exactly two output positions $n$ at which the error vanishes for all signals $u, v \in \mathbb{R}^N$:
\begin{equation}
    e[N-1] \;=\; 0 \qquad \text{and} \qquad e[N/2 - 1] \;=\; 0
\end{equation}
These are the only $n \in \{0,\ldots,N-1\}$ satisfying $(n-k)\bmod N = n \oplus k$ for all $k$; in particular, no output permutation can reconcile the two convolutions for all inputs.
\end{theorem}

\noindent\textit{Proof.}

\smallskip
\noindent\textit{(i) $e[N-1] = 0$.} Since $N-1=(1\cdots1)_2$, XOR with $N-1$ is bitwise complementation: $(N-1)\oplus k = N-1-k$, which equals $(N-1-k)\bmod N$ since $N-1-k\in\{0,\ldots,N-1\}$; thus every term of $y_{\mathrm{C}}[N-1]$ matches the corresponding term of $y_{\mathrm{D}}[N-1]$.

\smallskip
\noindent\textit{(ii) $e[N/2 - 1] = 0$.} Write $k = c_k\cdot N/2 + k'$ with $c_k\in\{0,1\}$ and $0\leq k'<N/2$. Since $n=(0\,\underbrace{1\cdots 1}_{m-1})_2$, XOR with $n$ inverts the lower $m-1$ bits and leaves the leading bit of $k$ unchanged. \textbf{Case} $c_k=0$: $n\oplus k = N/2-1-k'$ and $(n-k)\bmod N = N/2-1-k'$ (no modular reduction since $0\leq n-k < N/2$). \textbf{Case} $c_k=1$: $n\oplus k = N-1-k'$ and $(n-k)\bmod N = (-1-k')\bmod N = N-1-k'$ (since $-N<-1-k'<0$). Both cases agree.

\smallskip
\noindent\textit{(iii) Uniqueness.} We show no $n \in \{0,\ldots,N-2\}\setminus\{N/2-1\}$ satisfies the condition for all $k$. Use the witness $k = N-1$: since $(n-(N-1))\bmod N = n+1$ for $n < N-1$ (as $n-N+1<0$, adding $N$ gives $n+1$) and $n\oplus(N-1)=N-1-n$ (XOR with all-ones is bitwise complement), the condition requires $n+1 = N-1-n$, i.e.\ $n = N/2-1$. Any $n\in\{0,\ldots,N-2\}\setminus\{N/2-1\}$ therefore fails at $k=N-1$; the remaining case $n=N-1$ is already a zero-error position by part~(i). Together with parts (i) and (ii), the universal zero-error set is exactly $\{N/2-1,N-1\}$.

\smallskip
\noindent\textit{(iv) No-permutation claim.} Suppose some permutation $\phi$ satisfied $y_{\mathrm{D}}[\phi^{-1}(n)] = y_{\mathrm{C}}[n]$ for all $n$, $u$, $v$. The resulting bilinear-form equality requires $\phi^{-1}(n)\oplus k = (n-k)\bmod N$ for all $n$, $k$; setting $k=0$ gives $\phi^{-1}(n)=n$ for all $n$, so $\phi=\mathrm{id}$. But then the condition reduces to $(n-k)\bmod N = n\oplus k$ for all $n$, $k$, which part~(iii) shows fails for all $n\notin\{N/2-1,N-1\}$. $\square$

\medskip
\begin{corollary}[Dual zero-error input positions]\label{cor:dual-input}
The only $k \in \{0,\ldots,N-1\}$ satisfying $(n-k)\bmod N = n \oplus k$ for all $n$ are $k \in \{0,\,N/2\}$. In particular, columns $k=0$ and $k=N/2$ of $\mathbf{C}^u-\mathbf{H}^u$ are identically zero for all $u$.
\end{corollary}

\noindent\textit{Proof.} $k=0$ is immediate (since $(n-0)\bmod N = n = n\oplus 0$ for all $n$). For $k=N/2$: if $n \geq N/2$, then $(n-N/2)\bmod N = n-N/2$ (no wrap-around) and $n\oplus N/2 = n-N/2$ (clears the leading bit); if $n < N/2$, then $(n-N/2)\bmod N = n+N/2$ (wraps) and $n\oplus N/2 = n+N/2$ (sets the leading bit). Both cases agree. Uniqueness: for $k\neq 0$, at $n=0$, $(0-k)\bmod N = N-k$ and $0\oplus k = k$, so $N-k=k$ forces $k=N/2$. Since $k=0$ is already accounted for, these are the only two solutions. $\square$

\subsection{General Structure of the Error at an Arbitrary Position}
\label{sec:general-structure}

Having identified the positions where the error vanishes universally, we now derive a closed-form expression for $e[n]$ at an arbitrary output position $n$.

\noindent Define $\tilde{u}_n[j] \triangleq u[(n-j)\bmod N] - u[n \oplus j]$ for each output position $n$. Substituting $j=(n-k)\bmod N$ into $y_{\mathrm{C}}[n]$ and $j=n\oplus k$ into $y_{\mathrm{D}}[n]$ (both bijections on $\{0,\ldots,N-1\}$) and subtracting gives:
\begin{equation}
    \label{eq:general-closed-form}
    e[n] \;=\; \langle \tilde{u}_n,\, v \rangle \;=\; \sum_{j=0}^{N-1} \tilde{u}_n[j]\, v[j]
\end{equation}

\medskip
\begin{definition}[Mismatch permutation]\label{def:rho}
For $n \in \{0,\ldots,N-1\}$, the \emph{mismatch permutation} $\rho_n$ is:
\begin{equation}
    \label{eq:rho-def}
    \rho_n(j) \;=\; n \oplus \bigl((n-j)\bmod N\bigr)
\end{equation}
Since $\rho_n$ is the composition of $j\mapsto(n-j)\bmod N$ and $a\mapsto n\oplus a$, both bijections, it is a bijection on $\{0,\ldots,N-1\}$.
\end{definition}

\noindent The \emph{matching set} $S_n \triangleq \{j \in \{0,\ldots,N-1\} : (n-j)\bmod N = n \oplus j\}$ is the fixed-point set of $\rho_n$. Since $\tilde{u}_n[j] = 0$ whenever $j \in S_n$ (the circular and XOR indices agree there), the support of $\tilde{u}_n$ is contained in $S_n^{\mathsf{c}}$, with equality for generic $u$. Recall that $x$ is a \emph{submask} of $y$ if every set bit of $x$ is also set in $y$; the matching set has the following explicit form.

\medskip
\begin{theorem}[Matching set structure]\label{thm:matching-set}
For all $N = 2^m$ and $n \in \{0,\ldots,N-1\}$, with $n' \triangleq n \bmod (N/2)$:
\begin{equation}
    S_n \;=\; \bigl\{j \in \{0,\ldots,N-1\} \;:\; (j \bmod N/2) \text{ is a submask of } n'\bigr\}
\end{equation}
and $|S_n| = 2^{\,\mathrm{popcount}(n')\,+\,1}$, where $\mathrm{popcount}(n')$ denotes the number of $1$-bits in $n'$.
\end{theorem}

\medskip
\noindent\textit{Proof.} See Appendix~\ref{app:matching-set-proof}. $\square$

\begin{corollary}[Extremes of the matching-set spectrum]\label{cor:extremes}
$|S_n|$ is maximized at $n \in \{N/2-1,\,N-1\}$, giving $|S_n| = N$ (the universal zero-error output positions of Theorem~\ref{thm:zero-error}), and minimized at $n \in \{0,\,N/2\}$, giving $S_n = \{0,\,N/2\}$ (the universal zero-error input positions of Corollary~\ref{cor:dual-input}).
\end{corollary}

\noindent\textit{Proof.} Substitute $n' = N/2-1$ ($\mathrm{popcount}(n') = m-1$, giving $|S_n| = 2^m = N$) and $n' = 0$ ($\mathrm{popcount}(n') = 0$, giving $|S_n| = 2$) into Theorem~\ref{thm:matching-set}. $\square$

\subsection{Null-Space Structure of the Error Operator}
\label{sec:null-space}

The second result characterizes the rank and null-space structure of the error operator: the rank is generically $N-m-1$ (nearly full), and the null space has dimension $m+1$ (logarithmic in $N$). Define $\mathrm{lsb}(j) = p$ to mean $2^p \mid j$ and $2^{p+1} \nmid j$ (the position of the least significant $1$-bit of $j$). As shown in Appendix~\ref{app:rank-tightness}, the null-space dimension equals the number of orbits of $G = \langle\rho_n : n=0,\ldots,N-1\rangle$ acting on $\{0,\ldots,N-1\}$.


\begin{theorem}[Null-space structure]\label{thm:null-space}
For all $u \in \mathbb{R}^N$ with $N = 2^m$:
\begin{equation}
    \mathrm{rank}\bigl(\mathbf{C}^u - \mathbf{H}^u\bigr) \;\leq\; N - m - 1
\end{equation}
and $\dim\ker(\mathbf{C}^u - \mathbf{H}^u) \geq m+1$. The null space contains the $m+1$ linearly independent vectors
\begin{equation}
    \mathbf{b}_0 \quad\text{and}\quad \mathbf{f}_p \;\triangleq\; \sum_{\substack{j=1 \\ \mathrm{lsb}(j)=p}}^{N-1} \mathbf{b}_j, \quad p = 0, 1, \ldots, m-1
\end{equation}
Here $\mathbf{b}_j \in \mathbb{R}^N$ is the $j$-th standard basis vector. Moreover, for generic $u$ (all $u \in \mathbb{R}^N$ outside a measure-zero algebraic set):
\begin{equation}
    \mathrm{rank}\bigl(\mathbf{C}^u - \mathbf{H}^u\bigr) = N - m - 1
    \qquad\text{and}\qquad
    \ker(\mathbf{C}^u - \mathbf{H}^u) = \mathrm{span}\{\mathbf{b}_0,\,\mathbf{f}_0,\ldots,\mathbf{f}_{m-1}\}
\end{equation}
We call $\mathrm{span}\{\mathbf{b}_0,\mathbf{f}_0,\ldots,\mathbf{f}_{m-1}\}$ the \emph{universal zero-error subspace}.
\end{theorem}


\noindent\textit{Proof sketch.}

\smallskip
\noindent\textit{Linear independence and kernel membership of $\mathbf{b}_0$.} The lsb-level sets $\{0\}$ and $\{j:\mathrm{lsb}(j)=p\}$ for $p=0,\ldots,m-1$ partition $\{0,\ldots,N-1\}$ and are the supports of $\mathbf{b}_0,\mathbf{f}_0,\ldots,\mathbf{f}_{m-1}$; disjoint supports imply linear independence. Since $(n-0)\bmod N = n\oplus 0$ for all $n$, every row of $\mathbf{C}^u-\mathbf{H}^u$ vanishes on $\mathbf{b}_0$.

\smallskip
\noindent\textit{Kernel membership of $\mathbf{f}_p$.} Write $n=n_0+n_h\cdot 2^p$, $K=2^{m-p}$; row $n$ of $(\mathbf{C}^u-\mathbf{H}^u)\mathbf{f}_p$ equals $\sum_{\substack{d\,\mathrm{odd}\\1\le d<K}}[u[(n-d\cdot 2^p)\bmod N]-u[n\oplus d\cdot 2^p]]$. Both maps $d\mapsto(n_h-d)\bmod K$ and $d\mapsto n_h\oplus d$ send odd $d$ to values of parity opposite to $n_h$ (subtracting or XOR-ing an odd integer flips the low-order bit), are injective as restrictions of full-domain bijections, and hence are equal-image bijections onto the same $K/2$-element set; the $u$-index multisets coincide and the row equals zero. Full details in Appendix~\ref{app:fp-kernel-proof}.

\smallskip
\noindent\textit{Generic rank $N-m-1$.} It suffices to show $G=\langle\rho_n : n=0,\ldots,N-1\rangle$ acts on $\{0,\ldots,N-1\}$ with exactly $m+1$ orbits (Appendix~\ref{app:rank-tightness}): $\mathrm{lsb}$ is $G$-invariant (a bit-arithmetic case analysis on $\mathrm{lsb}(n)$ vs.\ $\mathrm{lsb}(j)$ shows $\rho_n$ preserves the least significant $1$-bit of $j$), each lsb-level forms a single orbit (step-$2$ adjacency on odd integers connects the entire level via a suitable witness choice of $n$), and $\{0\}$ is a fixed orbit since $\rho_n(0)=n\oplus n=0$. By dimension, $\ker(\mathbf{C}^u-\mathbf{H}^u)=\mathrm{span}\{\mathbf{b}_0,\mathbf{f}_0,\ldots,\mathbf{f}_{m-1}\}$ for generic $u$. Full proofs of kernel membership and rank tightness are in Appendices~\ref{app:fp-kernel-proof} and~\ref{app:rank-tightness}.

\medskip
\noindent\textbf{Remark (Non-generic filters).} The rank $N-m-1$ is tight for generic $u$, but can collapse to zero: any $u\in\mathrm{span}\{\mathbf{b}_0,\mathbf{f}_0,\ldots,\mathbf{f}_{m-1}\}$ satisfies $\mathbf{C}^u=\mathbf{H}^u$ and hence $\mathrm{rank}(\mathbf{C}^u-\mathbf{H}^u)=0$ (Appendix~\ref{app:rank-tightness}). Concretely, for $N=4$ the filter $u=[0,1,0,1]^\top$ assigns equal weight to the odd indices $\{1,3\}$ (the unique lsb-$0$ orbit) and zero elsewhere; direct computation confirms $\mathbf{C}^u=\mathbf{H}^u$.

\subsection{Expected Error}
\label{sec:expected-error}

The third result gives a closed-form expression for the expected error magnitude for random filters and signals, governed by a single alignment scalar. Central to our analysis is the \emph{alignment scalar}:
\begin{equation}
    \label{eq:Q-def}
    Q(u) \;\triangleq\; \sum_{n,k} u\bigl[(n-k)\bmod N\bigr]\,u[n\oplus k]
\end{equation}
which measures the index agreement between $\mathbf{C}^u$ and $\mathbf{H}^u$. Since $(n-k)\bmod N$ and $n\oplus k$ are each bijections in $k$ for fixed $n$, both matrices satisfy $\|\mathbf{C}^u\|_F^2 = \|\mathbf{H}^u\|_F^2 = N\|u\|_2^2$; the standard expansion $\|A-B\|_F^2 = \|A\|_F^2 - 2\langle A,B\rangle_F + \|B\|_F^2$ then gives:
\begin{equation}
    \label{eq:frobenius}
    \|\mathbf{C}^u-\mathbf{H}^u\|_F^2 \;=\; 2N\|u\|_2^2 - 2Q(u)
\end{equation}
For fixed $u$ and i.i.d.\ $v$ with zero mean and variance $\sigma_v^2$, expanding:
\begin{equation}
    \label{eq:Av-expand}
    \mathbb{E}_v\bigl[\|(\mathbf{C}^u-\mathbf{H}^u)v\|_2^2\bigr] \;=\; \sum_{n}\sum_{k,\ell} (\mathbf{C}^u-\mathbf{H}^u)_{nk}\,(\mathbf{C}^u-\mathbf{H}^u)_{n\ell}\,\mathbb{E}[v[k]v[\ell]] \;=\; \|\mathbf{C}^u-\mathbf{H}^u\|_F^2\,\sigma_v^2
\end{equation}
where cross-terms $k\neq\ell$ vanish because $v[k]$ are independent with zero mean. Substituting gives:
\begin{equation}
    \label{eq:fixed-filter-error}
    \mathbb{E}_v\bigl[\|\mathbf{e}\|_2^2\bigr] \;=\; \bigl(2N\|u\|_2^2 - 2Q(u)\bigr)\sigma_v^2
\end{equation}

\medskip
\begin{theorem}[Expected squared error, i.i.d.\ inputs]\label{thm:expected-error}
Let $u[0],\ldots,u[N-1]$ be i.i.d.\ with zero mean and variance $\sigma_u^2$, let $v[0],\ldots,v[N-1]$ be i.i.d.\ with zero mean and variance $\sigma_v^2$, and let $u$ and $v$ be independent of each other. Then:
\begin{equation}
    \label{eq:expected-error}
    \mathbb{E}\bigl[\|\mathbf{e}\|_2^2\bigr] \;=\; 2\sigma_u^2\sigma_v^2\bigl(N^2 - 4\cdot 3^{m-1}\bigr)
\end{equation}
\end{theorem}

\noindent\textit{Proof.} Using equation~\eqref{eq:frobenius}, taking expectation over $u$: $\mathbb{E}[\|u\|_2^2] = N\sigma_u^2$, so $\mathbb{E}[2N\|u\|_2^2] = 2N^2\sigma_u^2$. For the cross-term, $\mathbb{E}[Q(u)] = M\sigma_u^2$ where $M = \sum_{n,k}\mathbb{1}[(n-k)\bmod N = n\oplus k] = \sum_{n=0}^{N-1}|S_n|$ counts matching index pairs: for a matching pair both terms index the same entry of $u$, giving $\mathbb{E}[u[(n-k)\bmod N]\,u[n\oplus k]] = \sigma_u^2$; for a non-matching pair the indices are distinct, so $\mathbb{E}[u[a]u[b]] = 0$ by independence.

By Theorem~\ref{thm:matching-set}, $|S_n| = 2^{\mathrm{popcount}(n')+1} = 2\cdot 2^{\mathrm{popcount}(n')}$, so:
\begin{align}
    M &= \sum_{n=0}^{N-1}|S_n|
      = 2\sum_{n=0}^{N-1}2^{\mathrm{popcount}(n\bmod N/2)} \notag\\
      &= 4\sum_{\nu=0}^{N/2-1}2^{\mathrm{popcount}(\nu)}
      = 4\cdot 3^{m-1}
\end{align}
where the third equality substitutes $\nu = n\bmod(N/2)$, folding each $\nu \in \{0,\ldots,N/2-1\}$ into its two preimages $n=\nu$ and $n=\nu+N/2$, and the fourth uses $\sum_{\nu=0}^{N/2-1}2^{\mathrm{popcount}(\nu)} = 3^{m-1}$ (the $m-1$ bit positions of $\nu$ each contribute a factor $1+2=3$). Substituting $M = 4\cdot 3^{m-1}$ into $\mathbb{E}[\|\mathbf{C}^u-\mathbf{H}^u\|_F^2] = 2\sigma_u^2(N^2-M)$ gives $\mathbb{E}[\|\mathbf{C}^u-\mathbf{H}^u\|_F^2] = 2\sigma_u^2(N^2 - 4\cdot 3^{m-1})$. Conditioning on $u$ and applying $\mathbb{E}_v[\|(\mathbf{C}^u-\mathbf{H}^u)v\|_2^2] = \|\mathbf{C}^u-\mathbf{H}^u\|_F^2\sigma_v^2$ yields equation~\eqref{eq:expected-error}. $\square$

\medskip
\begin{corollary}[Error-to-output energy ratio]\label{cor:ratio}
Under the i.i.d.\ model of Theorem~\ref{thm:expected-error},
\begin{equation}
    \label{eq:ratio}
    \frac{\mathbb{E}[\|\mathbf{e}\|_2^2]}{\mathbb{E}[\|y_{\mathrm{C}}\|_2^2]} \;=\; 2\!\left(1-\!\left(\tfrac{3}{4}\right)^{\!m-1}\right)
\end{equation}
which approaches $2$ as $N\to\infty$.
\end{corollary}

\noindent\textit{Proof.} The denominator follows by applying the same expansion as equation~\eqref{eq:Av-expand} to $\mathbf{C}^u$ alone: $\mathbb{E}_v[\|\mathbf{C}^u v\|_2^2] = \|\mathbf{C}^u\|_F^2\sigma_v^2 = N\|u\|_2^2\sigma_v^2$, and averaging over $u$ gives $\mathbb{E}[\|y_{\mathrm{C}}\|_2^2] = N\mathbb{E}[\|u\|_2^2]\sigma_v^2 = N^2\sigma_u^2\sigma_v^2$. Dividing equation~\eqref{eq:expected-error} by this and simplifying yields equation~\eqref{eq:ratio}. $\square$

\medskip
\noindent\textbf{Remark (Error energy exceeds output energy asymptotically).} For large $N$ the substitution error carries twice the energy of the intended output. The mechanism is decorrelation: $\mathbb{E}[\|\mathbf{e}\|_2^2] = 2\,\mathbb{E}[\|y_{\mathrm{C}}\|_2^2] - 2\,\mathbb{E}[\langle y_{\mathrm{C}}, y_{\mathrm{D}}\rangle]$, and the normalized cross-term $(3/4)^{m-1}$ decays to zero as $N\to\infty$ (Theorem~\ref{thm:expected-error}), so the two energies accumulate additively rather than cancelling.

\medskip
\noindent The alignment scalar $Q(u)$ is itself a quadratic form whose $N$-eigenspace is exactly the universal zero-error subspace of Section~\ref{sec:null-space}, tying the rank and alignment results together.

\medskip
\begin{theorem}[Range and eigenspace of the alignment scalar $Q(u)$]\label{thm:lambda-max}
For all $N = 2^m$ and $u \in \mathbb{R}^N$:
\begin{equation}
    -N\|u\|_2^2 \;\leq\; Q(u) \;\leq\; N\,\|u\|_2^2
\end{equation}
The upper bound is achieved exactly on the universal zero-error subspace $\mathrm{span}\{\mathbf{b}_0, \mathbf{f}_0, \ldots, \mathbf{f}_{m-1}\}$ of Theorem~\ref{thm:null-space}.
\end{theorem}

\noindent\textit{Proof.} Both bounds follow from Cauchy--Schwarz: $|Q(u)| = |\langle\mathbf{C}^u,\mathbf{H}^u\rangle_F| \leq \|\mathbf{C}^u\|_F\|\mathbf{H}^u\|_F = N\|u\|_2^2$. Upper bound and eigenspace characterization: Appendix~\ref{app:quadratic-form}. $\square$

\section{Discussion}

\paragraph{Matching-set geometry as the unifying structure.}
The fixed‑point structure of $\rho_n$ underlies all three classes of results: the matching set determines the zero‑error conditions (Sections~\ref{sec:foundational}–\ref{sec:general-structure}), the way $\rho_n$ partitions indices into orbits determines the null‑space dimension (Section~\ref{sec:null-space}), and the total number of matching pairs determines the closed‑form expression for the expected error (Section~\ref{sec:expected-error}). The two universally zero‑error output positions (Theorem~\ref{thm:zero-error}) and the two universally zero‑error input positions $k \in \{0,\,N/2\}$ (Corollary~\ref{cor:dual-input}) are those where every index matches. Furthermore, no output permutation can reconcile the two convolutions for arbitrary inputs (Theorem~\ref{thm:zero-error}). At all other positions, the error is a dot product between the input and a filter‑difference signal whose support equals the set of non‑matching indices; smaller support leads to smaller error.

\paragraph{Near-full rank and practical implications.}
The substitution error affects nearly all inputs: for generic filters, the error operator has rank $N-m-1$, leaving a zero‑error subspace of only $\log_2 N + 1$ dimensions. At $N=1024$, the zero-error subspace spans just $11$ of $1024$ dimensions ($1.1\%$); at $N=65536$, $17$ of $65536$ ($0.026\%$), shrinking rapidly with signal length. Hadamard layers proposed as drop-in replacements for DFT-based convolution~\cite{pan2021fast, pan2022block, pan2023htlayer} therefore compute a structurally distinct operation on all but a vanishingly small fraction of inputs. Combined with Corollary~\ref{cor:ratio}, this implies that the substitution error carries substantial energy for essentially all inputs. The normalized alignment $Q(u)/\|u\|_2^2$ (Section~\ref{sec:expected-error}) provides a single-number diagnostic: the closer this ratio is to its maximum value of $N$, the smaller the substitution error.

\paragraph{Filter alignment as the governing design variable.}
For random filters and signals, the error‑to‑output energy ratio approaches $2$ as signal length increases (Corollary~\ref{cor:ratio}): absent structural alignment, the substitution error asymptotically doubles the output energy. Filters $u$ lying in the zero-error subspace achieve maximum alignment ($Q(u)/\|u\|_2^2 = N$, the largest value permitted by Theorem~\ref{thm:lambda-max}) and incur no error, while the fraction of maximum alignment for generic i.i.d.\ filters satisfies $\mathbb{E}[Q(u)]/(N\cdot\mathbb{E}[\|u\|_2^2]) = \tfrac{4}{3}(3/4)^m \to 0$ as $N\to\infty$, placing them increasingly far from the zero-error regime. For a learned filter, the normalized alignment $Q(u)/\|u\|_2^2$ is bounded by $[-N,\,N]$ (Theorem~\ref{thm:lambda-max}), and whether training drives filters toward or away from the zero-error subspace is an open empirical question.


\section{Acknowledgements}

The authors thank Zahra Ronaghi and Saee Paliwal for supporting this project. The authors also thank Brian L. Evans of The University of Texas at Austin and Dilip Sarwate of the University of Illinois Urbana-Champaign for discussions that motivated this work. The authors declare no financial interest or conflicts.

\bibliographystyle{plainnat}
\bibliography{wht_conv}

\newpage

\setcounter{theorem}{0}
\appendix

\section{Proof of Theorem~\ref{thm:matching-set} (Matching Set Structure)}
\label{app:matching-set-proof}

\noindent\textit{Proof.} Write $n = n' + c_n \cdot N/2$ and $j = j' + c_j \cdot N/2$ with $n', j' \in \{0,\ldots,N/2-1\}$ and $c_n, c_j \in \{0,1\}$. In both cases ($c_n = c_j$ and $c_n \neq c_j$), the leading-bit contributions cancel: $n \oplus j = (n' \oplus j') + (c_n \oplus c_j)\cdot N/2$, and $(n-j)\bmod N$ evaluates as follows.

\medskip
\noindent\textbf{Case A} ($c_n = c_j$): $n - j = n' - j'$. For $j' \leq n'$: $(n-j)\bmod N = n'-j'$ (non-negative, less than $N/2 \leq N$, no reduction needed), and $n \oplus j = n' \oplus j'$ (leading bits cancel). For $j' > n'$: $(n-j)\bmod N = N + (n'-j') \geq N/2+1$, so the leading bit is $1$; but $n \oplus j = n' \oplus j' < N/2$ has leading bit $0$. No agreement.

\medskip
\noindent\textbf{Case B} ($c_n \neq c_j$): Sub-case $c_n=1, c_j=0$: $n-j = N/2 + n' - j' \in [1,N-1]$, no reduction is needed. Sub-case $c_n=0, c_j=1$: $n-j = n' - N/2 - j' < 0$, adding $N$ gives $N/2 + n' - j'$. In both sub-cases: $(n-j)\bmod N = N/2 + (n'-j')$ and $n \oplus j = (n' \oplus j') + N/2$ (since $c_n \oplus c_j = 1$). These agree iff $n'-j' = n' \oplus j'$.

\medskip
In all cases, the condition reduces to $n' - j' = n' \oplus j'$ with $j' \leq n'$. The identity $a - b = a \oplus b$ holds for non-negative integers iff $b$ is a submask of $a$ (the subtraction has no borrows). Therefore $(n-j)\bmod N = n \oplus j$ if and only if $j' \subseteq n'$.

The number of $j \in \{0,\ldots,N-1\}$ with $j' = j \bmod N/2 \subseteq n'$ is: $2^{\mathrm{popcount}(n')}$ choices for $j'$ (one per submask of $n'$) times $2$ choices for $c_j$, giving $|S_n| = 2^{\mathrm{popcount}(n')+1}$. $\square$

\section{Proof of Kernel Membership of $\mathbf{f}_p$ (Theorem~\ref{thm:null-space})}
\label{app:fp-kernel-proof}

\noindent\textit{Proof.} Row $n$ of $(\mathbf{C}^u-\mathbf{H}^u)\mathbf{f}_p$ equals:
\[
    \sum_{\substack{d \text{ odd} \\ 1 \leq d < 2^{m-p}}} \Bigl[u\bigl[(n - d\cdot 2^p)\bmod N\bigr] - u\bigl[n \oplus d\cdot 2^p\bigr]\Bigr]
\]
Write $n = n_0 + n_h\cdot 2^p$ with $n_0 = n\bmod 2^p$ and $n_h = \lfloor n/2^p\rfloor$. Since $d\cdot 2^p$ is a multiple of $2^p$, the low $p$ bits are unaffected by either operation:
\begin{align*}
    (n - d\cdot 2^p)\bmod N &= n_0 + \bigl((n_h - d)\bmod 2^{m-p}\bigr)\cdot 2^p \\
    n \oplus d\cdot 2^p &= n_0 + (n_h \oplus d)\cdot 2^p
\end{align*}
Let $K = 2^{m-p}$ and $P = \{\ell \in \{0,\ldots,K-1\} : \ell \not\equiv n_h \pmod{2}\}$; note $|P| = K/2$. As $d$ ranges over the $K/2$ odd values in $\{1,\ldots,K-1\}$:
\begin{itemize}
\item \emph{Image in $P$.} Subtracting (resp.\ XOR-ing) an odd $d$ flips the low-order bit of $n_h$, so $(n_h - d)\bmod K$ and $n_h \oplus d$ both have parity opposite to $n_h$; hence both maps send odd $d$ into $P$.
\item \emph{Injectivity on odd $d$.} The map $d\mapsto(n_h-d)\bmod K$ is a bijection on all of $\{0,\ldots,K-1\}$; a bijection on the full domain restricts to an injection on any subset. Similarly, $d\mapsto n_h\oplus d$ is a bijection on $\{0,\ldots,K-1\}$, hence injective on odd $d$.
\item \emph{Bijection onto $P$.} Each restricted map is an injection from a $K/2$-element domain into a $K/2$-element codomain $P$, so both restrictions are bijections onto $P$.
\end{itemize}
Since both maps have the same image $P$, the multisets $\{(n_h-d)\bmod K : d\text{ odd}\}$ and $\{n_h\oplus d : d\text{ odd}\}$ are equal. The corresponding $u$-indices $n_0 + \ell\cdot 2^p$ therefore coincide as multisets, the sums over $u[\cdot]$ are equal, and the difference is zero. $\square$

\section{Proof of Rank Tightness for Generic $u$ (Theorem~\ref{thm:null-space})}
\label{app:rank-tightness}

\noindent\textit{Proof.} We work over $\mathbb{Q}(u_0,\ldots,u_{N-1})$, treating the filter entries as algebraically independent indeterminates; the Conclusion converts this to the measure-theoretic statement.

\medskip
\noindent\textit{Step 1: Generic null space = functions constant on orbits.} Over $\mathbb{Q}(u_0,\ldots,u_{N-1})$, $(\mathbf{C}^u - \mathbf{H}^u)\mathbf{w} = 0$ requires:
\[
\sum_a u[a]\bigl(\mathbf{w}[(n-a)\bmod N] - \mathbf{w}[n\oplus a]\bigr) = 0 \quad\text{for all }n
\]
Let $q \in \mathbb{Q}[u_0,\ldots,u_{N-1}]$ be a common denominator for the entries $\mathbf{w}[j] \in \mathbb{Q}(u_0,\ldots,u_{N-1})$, and write $\mathbf{w}[j] = g[j]/q$ with $g[j] \in \mathbb{Q}[u_0,\ldots,u_{N-1}]$. Multiplying through by $q$ gives $\sum_a u[a](g[(n-a)\bmod N] - g[n\oplus a]) = 0$ as a polynomial identity; since $u_0,\ldots,u_{N-1}$ are algebraically independent, each coefficient vanishes: $g[(n-a)\bmod N] = g[n\oplus a]$ for all $n,a$, hence $\mathbf{w}[(n-a)\bmod N] = \mathbf{w}[n\oplus a]$ for all $n,a$. Substituting $b = (n-a)\bmod N$ (so $n\oplus a = \rho_n(b)$), this reads $\mathbf{w}[b] = \mathbf{w}[\rho_n(b)]$ for all $n,b$. That is, the null space over $\mathbb{Q}(u_0,\ldots,u_{N-1})$ equals the space of vectors constant on orbits of $G = \langle\rho_n : n = 0,\ldots,N-1\rangle$.

\medskip
\noindent\textit{Step 2: Orbit Lemma.} $G$ acts on $\{0,\ldots,N-1\}$ with exactly $m+1$ orbits: $\{0\}$ and $\{j : \mathrm{lsb}(j) = p\}$ for $p = 0,\ldots,m-1$. First, $\rho_n(0) = n \oplus ((n-0)\bmod N) = n \oplus n = 0$ for all $n$, so $\{0\}$ is fixed by every generator and hence forms a $G$-orbit on its own.

\smallskip
\emph{(a) lsb is $G$-invariant.} We show $\mathrm{lsb}(\rho_n(j)) = \mathrm{lsb}(j)$ for all $n,j$. Set $s = \mathrm{lsb}(n)$, $t = \mathrm{lsb}(j)$, write $n = 2^s \bar{n}$, $j = 2^t \bar{j}$ with $\bar{n},\bar{j}$ odd (here $\bar{\cdot}$ denotes the odd part after factoring out the power of $2$; distinct from $n' = n\bmod N/2$ used in the main body and Appendix~\ref{app:matching-set-proof}), and let $\alpha = (n-j)\bmod N$.

\begin{itemize}
\item $s > t$: Then $\alpha = 2^t \cdot (\text{odd})$, so bit $t$ of $\alpha$ is $1$. Since $s > t$, bit $t$ of $n$ is $0$. Thus bit $t$ of $n\oplus \alpha = 0\oplus 1 = 1$, and bits $0,\ldots,t-1$ of $n\oplus \alpha$ are zero. Hence $\mathrm{lsb}(\rho_n(j)) = t$.

\item $s < t$: Write $n - j = 2^s(\bar{n} - 2^{t-s}\bar{j})$. Since $\bar{n}$ is odd and $2^{t-s}\bar{j}$ is even, the factor $\bar{n} - 2^{t-s}\bar{j}$ is odd; set $c = (\bar{n} - 2^{t-s}\bar{j})\bmod 2^{m-s}$ (odd), so $\alpha = 2^s c$. Since both $n$ and $\alpha$ are divisible by $2^s$, we have $n\oplus \alpha = 2^s(\bar{n}\oplus c)$, so $\mathrm{lsb}(\rho_n(j)) = s + \mathrm{lsb}(\bar{n}\oplus c)$. Bits $0,\ldots,t-s-1$ of $c$ equal those of $\bar{n}$: since $2^{t-s}\bar{j}$ is zero in those bit positions and binary subtraction carries propagate toward higher-order bits, zeros in positions $0,\ldots,t-s-1$ of $2^{t-s}\bar{j}$ guarantee those bit positions of $c$ match those of $\bar{n}$. Hence bits $0,\ldots,t-s-1$ of $\bar{n}\oplus c$ are zero. At bit $t-s$: since $\bar{j}$ is odd, bit $t-s$ of $2^{t-s}\bar{j}$ is $1$, and there is no borrow from below, so subtracting flips bit $t-s$ of $\bar{n}$; hence the $(t-s)$-th bit of $\bar{n}\oplus c$ is $1$. Therefore $\mathrm{lsb}(\bar{n}\oplus c) = t-s$ and $\mathrm{lsb}(\rho_n(j)) = s+(t-s) = t$.

\item $s = t$: Then $\alpha = 2^t(\text{even})$, so $\mathrm{lsb}(\alpha) > t$, bit $t$ of $\alpha = 0$, bit $t$ of $n = 1$. Thus, bit $t$ of $n\oplus \alpha = 1$ and bits $0,\ldots,t-1$ are zero. Hence $\mathrm{lsb}(\rho_n(j)) = t$.
\end{itemize}

\smallskip
\emph{Carry-theoretic reading.} The case analysis has a compact restatement: the first borrow in the subtraction $n-j$ occurs at bit $t=\mathrm{lsb}(j)$, since bits $0,\ldots,t-1$ of $j$ are zero and bit $t$ of $j$ is $1$. Bits $0,\ldots,t-1$ of $\alpha=(n-j)\bmod N$ therefore agree with those of $n$, and bit $t$ differs, giving $\mathrm{lsb}(n\oplus \alpha)=t$ directly. Correspondingly, $j\in S_n$ iff the subtraction of their lower $m-1$ bits is borrow-free---equivalently, iff $(j\bmod N/2)$ is a submask of $(n\bmod N/2)$ (Theorem~\ref{thm:matching-set}).

\smallskip
\emph{(b) Transitivity.} We show $G$ acts transitively on $\{j : \mathrm{lsb}(j) = p\}$ for each $p$. Writing $n = 2^p n_h + n_0$ with $n_0 = n\bmod 2^p$, the action of $\rho_n$ on $2^p\tilde{j}$ (odd $\tilde{j}$) is:
\[
\rho_n(2^p \tilde{j}) = 2^p\bigl(n_h \oplus ((n_h - \tilde{j})\bmod K)\bigr) = 2^p\,\rho^{(K)}_{n_h}(\tilde{j})
\]
where $\rho^{(K)}_a(b) = a\oplus((a-b)\bmod K)$ and $K = 2^{m-p}$. Since $d\cdot 2^p$ is divisible by $2^p$, the low $p$ bits $n_0$ are unaffected by both subtraction and XOR, and cancel from $\rho_n(2^p\tilde{j})$. It suffices to show $G^{(K)} = \langle\rho^{(K)}_{n_h}\rangle$ acts transitively on the set $O_K$ of odd elements of $\{0,\ldots,K-1\}$. For $K=2$ (i.e., $p=m-1$), $O_K = \{1\}$ is a single element and transitivity is immediate. For $K \geq 4$, for any odd $\tilde{j} \leq K-3$, choose $n_h = \tilde{j}+1$ (even, so $n_h - \tilde{j} = 1$):
\[
\rho^{(K)}_{\tilde{j}+1}(\tilde{j}) \;=\; (\tilde{j}+1)\oplus 1 \;=\; \tilde{j}+2
\]
since $\tilde{j}+1$ is even (bit 0 is $0$), so XOR with $1$ sets bit 0, giving $\tilde{j}+2$. Thus, for every odd $\tilde{j} \leq K-3$, the elements $\tilde{j}$ and $\tilde{j}+2$ lie in the same $G^{(K)}$-orbit. Since $O_K = \{1,3,5,\ldots,K-1\}$ and the step-$2$ adjacency connects $1\sim 3\sim 5\sim\cdots\sim K-1$, transitivity of the orbit equivalence relation implies all elements of $O_K$ share a single orbit. Hence $G^{(K)}$ acts transitively on $O_K$.

\medskip
\noindent\textit{Conclusion.} Steps~1 and~2 together show that the null space of $\mathbf{C}^u-\mathbf{H}^u$ over $\mathbb{Q}(u_0,\ldots,u_{N-1})$ is exactly the space of functions constant on the $m+1$ orbits of $G$, which has dimension $m+1$; hence $\mathrm{rank}(\mathbf{C}^u-\mathbf{H}^u) = N-(m+1) = N-m-1$ over $\mathbb{Q}(u_0,\ldots,u_{N-1})$. By definition of rank over a field, some $(N-m-1)\times(N-m-1)$ submatrix has nonzero determinant in $\mathbb{Q}(u_0,\ldots,u_{N-1})$; since every entry of $\mathbf{C}^u-\mathbf{H}^u$ is linear in $u_0,\ldots,u_{N-1}$, this determinant is a nonzero polynomial $P\in\mathbb{Z}[u_0,\ldots,u_{N-1}]$ (not merely a nonzero rational function). Since the entries of $\mathbf{C}^u-\mathbf{H}^u$ are linear in $u$, every minor is a polynomial, so $\{u:\mathrm{rank}\geq N-m-1\}$ is Zariski open. A nonzero polynomial over $\mathbb{R}$ cannot vanish on all of $\mathbb{R}^N$, so $P$ takes a nonzero value at some $u\in\mathbb{R}^N$ and the Zariski-open set is nonempty. Its complement is therefore a proper Zariski-closed set in $\mathbb{R}^N$, which has Lebesgue measure zero, so $\mathrm{rank}(\mathbf{C}^u-\mathbf{H}^u) = N-m-1$ for generic $u\in\mathbb{R}^N$. Finally, by Theorem~\ref{thm:null-space}, $\mathrm{span}\{\mathbf{b}_0, \mathbf{f}_0, \ldots, \mathbf{f}_{m-1}\} \subseteq \ker(\mathbf{C}^u - \mathbf{H}^u)$ for every $u$; for generic $u$, $\dim\ker = m+1 = \dim\,\mathrm{span}\{\mathbf{b}_0, \mathbf{f}_0, \ldots, \mathbf{f}_{m-1}\}$, so the inclusion is an equality. $\square$

\medskip

\noindent\textbf{Remark (Numerical illustration; not part of the proof).} The computations below verify the rank result of Theorem~\ref{thm:null-space} for the witness filter $u^* = [0,1,\ldots,N-1]^\top$, for which $(\mathbf{C}^{u^*}-\mathbf{H}^{u^*})_{n,k} = (n-k)\bmod N - (n\oplus k)$; they are provided for illustrative purposes only.


\smallskip
\noindent\textit{Case $N=4$, $m=2$.} The matrix $\mathbf{C}^{u^*}-\mathbf{H}^{u^*}$ is:
\[
\begin{pmatrix}
 0 &  2 & 0 & -2 \\
 0 &  0 & 0 &  0 \\
 0 & -2 & 0 &  2 \\
 0 &  0 & 0 &  0
\end{pmatrix}
\]
Rows $n=1$ and $n=3$ are zero (the universal zero-error output positions $N/2-1$ and $N-1$); row~$0$ and row~$2$ sum to zero. Hence $\mathrm{rank}=1=N-m-1$, and the null space is $\mathrm{span}\{\mathbf{b}_0,\mathbf{f}_0,\mathbf{f}_1\}$ where $\mathbf{b}_0=[1,0,0,0]^\top$, $\mathbf{f}_0=[0,1,0,1]^\top$ (odd indices), $\mathbf{f}_1=[0,0,1,0]^\top$ (index~$2$).

\smallskip
\noindent\textit{Case $N=8$, $m=3$.} The eight rows of $\mathbf{C}^{u^*}-\mathbf{H}^{u^*}$ are:
\begin{align*}
r_0 &= [\phantom{-}0,\phantom{-}6,\phantom{-}4,\phantom{-}2,\phantom{-}0,-2,-4,-6],\\
r_1 &= [\phantom{-}0,\phantom{-}0,\phantom{-}4,\phantom{-}4,\phantom{-}0,\phantom{-}0,-4,-4],\\
r_2 &= [\phantom{-}0,-2,\phantom{-}0,\phantom{-}6,\phantom{-}0,-2,\phantom{-}0,-2],\\
r_3 &= [\phantom{-}0,\phantom{-}0,\phantom{-}0,\phantom{-}0,\phantom{-}0,\phantom{-}0,\phantom{-}0,\phantom{-}0],\\
r_4 &= [\phantom{-}0,-2,-4,-6,\phantom{-}0,\phantom{-}6,\phantom{-}4,\phantom{-}2],\\
r_5 &= [\phantom{-}0,\phantom{-}0,-4,-4,\phantom{-}0,\phantom{-}0,\phantom{-}4,\phantom{-}4],\\
r_6 &= [\phantom{-}0,-2,\phantom{-}0,-2,\phantom{-}0,-2,\phantom{-}0,\phantom{-}6],\\
r_7 &= [\phantom{-}0,\phantom{-}0,\phantom{-}0,\phantom{-}0,\phantom{-}0,\phantom{-}0,\phantom{-}0,\phantom{-}0].
\end{align*}
Rows $r_3$ and $r_7$ are zero (the universal zero-error positions $n=N/2-1$ and $n=N-1$). The six nonzero rows satisfy exactly two linear dependencies: $r_1+r_5=0$ and $r_0+r_2+r_4+r_6=0$. Hence $\mathrm{rank}=6-2=4=N-m-1$, and the null space is $\mathrm{span}\{\mathbf{b}_0,\mathbf{f}_0,\mathbf{f}_1,\mathbf{f}_2\}$ where $\mathbf{b}_0=\mathbf{e}_0$, $\mathbf{f}_0=\mathbf{e}_1+\mathbf{e}_3+\mathbf{e}_5+\mathbf{e}_7$ (odd indices), $\mathbf{f}_1=\mathbf{e}_2+\mathbf{e}_6$ (lsb-$1$ indices), $\mathbf{f}_2=\mathbf{e}_4$ (lsb-$2$ index).

\section{Quadratic Form Analysis: Proof of Theorem~\ref{thm:lambda-max}}
\label{app:quadratic-form}


\begin{thebibliography}{99}
\bibitem[]{PueschelMoura2008_ASP_1DSpace} P{\"u}schel, Markus; Moura, Jos{\'e} M. F.. Algebraic Signal Processing Theory: 1-D Space. IEEE Transactions on Signal Processing 56 (2008) 3586--3599
\bibitem[]{PueschelMoura2008_ASP_Foundation1DTime} P{\"u}schel, Markus; Moura, Jos{\'e} M. F.. Algebraic Signal Processing Theory: Foundation and 1-D Time. IEEE Transactions on Signal Processing 56 (2008) 3572--3585
\bibitem[]{Terras_1999} Terras, Audrey. Fourier Analysis on Finite Groups and Applications. (1999)
\bibitem[]{ahmed1975orthogonal} Nasir Ahmed; Kamisetty R. Rao. Orthogonal Transforms for Digital Signal Processing. (1975)
\bibitem[]{beauchamp1975walsh} Beauchamp, K. G.. Walsh Functions and Their Applications. (1975)
\bibitem[]{doi:10.1049/el:19730172} M. N. Gulamhusein. Simple Matrix-Theory Proof of the Discrete Dyadic Convolution Theorem. Electronics Letters 9 (1973) 238-239
\bibitem[]{fino1976unified} Fino, Bernard J.; Algazi, V. Ralph. A Unified Treatment of Discrete Fast Unitary Transforms. SIAM Journal on Computing 6 (1977) 700--717
\bibitem[]{pan2021fast} Pan, Hongyi; Badawi, Diaa; Cetin, Ahmet Enis. Fast Walsh-Hadamard Transform and Smooth-Thresholding Based Binary Layers in Deep Neural Networks. Proceedings of the IEEE/CVF Conference on Computer Vision and Pattern Recognition (CVPR) Workshops (2021) 4650-4659
\bibitem[]{pan2022block} Pan, Hongyi; Badawi, Diaa; Cetin, Ahmet Enis. Block Walsh-Hadamard Transform Based Binary Layers in Deep Neural Networks. ACM Transactions on Embedded Computing Systems 21 (2022)
\bibitem[]{pan2023htlayer} Pan, Hongyi; Zhu, Xin; Atici, Salih; Cetin, Ahmet Enis. A Hybrid Quantum-Classical Approach based on the Hadamard Transform for the Convolutional Layer. Proceedings of the 40th International Conference on Machine Learning (2023)
\bibitem[]{shanks1969ITCmp.100.457S} Shanks, J.~L.. Computation of the Fast Walsh-Fourier Transform. IEEE Transactions on Computers 18 (1969) 457-459
\bibitem[]{uvsakova2002walsh} Usakova, Anna; Kotuliaková, Jana; Zajac, Michal. Walsh--Hadamard Transformation of a Convolution. Radioengineering 11 (2002) 40--42
\bibitem[]{walsh1923closed} Walsh, Joseph L.. A Closed Set of Normal Orthogonal Functions. American Journal of Mathematics 45 (1923) 5--24
\bibitem[]{yarlagadda1997hadamard} Yarlagadda, Rao K.; Hershey, John E.. Hadamard Matrix Analysis and Synthesis: With Applications to Communications and Signal/Image Processing. (1997)
\end{thebibliography}
\end{document}
